\section{极化传播子与两点以上格林函数}

Gross 第 16、27 章指出：势能与密度涨落由极化传播子 $\Pi$ 决定；非局域单体算符的关联则需要两点格林函数 $G_2$。运动方程把它们连成层级。

\subsection{极化传播子}

\begin{equation}
  i\Pi(x t,x't')
  =\bigl\langle
  T\bigl[\delta\hat n(xt)\,\delta\hat n(x't')\bigr]
  \bigr\rangle,
  \qquad
  \delta\hat n=\hat n-\langle n\rangle.
\label{eq:Pi_def}
\end{equation}
相互作用能中与涨落有关的部分为
\begin{equation}
  \langle V\rangle
  =\frac12\int\mathrm{d}x\,\mathrm{d}x'\,
  v(x,x')
  \bigl[
  i\Pi(x0,x'0)+\langle n(x)\rangle\langle n(x')\rangle
  -\delta(x-x')\langle n(x)\rangle
  \bigr].
\end{equation}
（最后一项扣除自作用。）线性密度响应 $\chi_{nn}$ 即 $\Pi$ 的 retarded 版，与第 5 章 Kubo 公式同一对象。

\subsection{两点格林函数}

\begin{equation}
  i^2 G_2(1,2,3,4)
  =\bigl\langle
  T\,\psi(1)\psi(2)\psi^\dagger(3)\psi^\dagger(4)
  \bigr\rangle.
\label{eq:G2}
\end{equation}
费米反对称：
\begin{equation}
  G_2(1,2,3,4)
  =-G_2(2,1,3,4)
  =-G_2(1,2,4,3)
  =G_2(2,1,4,3).
\end{equation}
四个时间有 $24$ 种顺序，反对称后独立情况减为 $6$ 类（Gross 第 16.4 节）。$G_2$ 足够计算任意两个单体算符的关联 $\langle AB\rangle-\langle A\rangle\langle B\rangle$。

\subsection{与 $G$ 的运动方程}

对 $i\partial_{t_1}G(1,1')$ 使用 $[H,\psi]$，相互作用产生
\begin{equation}
  \bigl(i\partial_{t_1}-h_0(1)\bigr)G(1,1')
  =\delta(1,1')
  -i\int\mathrm{d}2\,v(1,2)G_2(1,2,1',2^+).
\end{equation}
再对 $G_2$ 写方程出现 $G_3$，即 Martin--Schwinger 层级。截断
$G_2\approx G G - G G$（交换）给出 HF；把 $G_2$ 用不可约顶点 $\Gamma$ 参数化则给出 Bethe--Salpeter 方程（见后文）。

\begin{note}
RPA 相当于在 $\Pi$ 通道只保留气泡：$\Pi^{-1}=\Pi_0^{-1}-v$。这是“高级视角下的等离激元”（Gross 第 28 章）：集体极点来自有效相互作用的零点，而非单粒子 $G$ 的极点。
\end{note}
